\chapter{고찰}
\label{ch:discussion}

\section{정보 입도 한계}
\label{sec:disc-granularity}

제 \ref{ch:results} 장의 결과는 하나의 해석으로 수렴한다. 분 단위 공개 텔레메트리는 교전 결과의 \emph{거시적} 결정 요인(누적 자원 격차, 오브젝트 통제, 생존 인원)을 담고 있지만 \emph{미시적} 결정 요인(스킬 적중, 진입 타이밍, 쿨다운)은 담고 있지 않다. 거시 요인이 결과를 지배하는 후반·일방 교전에서는 AUC가 $.8$에 이르고, 미시 요인이 결과를 가르는 대등한 교전에서는 우연 수준에 가깝다. 이는 관측 정보의 상호정보량 관점에서 자연스럽다. 라벨 $Y$와 관측 $X$의 상호정보량 $I(X;Y)$는 관측 해상도가 거칠수록 작아지며, 어떤 학습기도 $I(X;Y)$가 허용하는 이상의 판별력을 얻을 수 없다. 리플레이 기반의 초 단위 관측을 전제한 선행 연구 \cite{katona2019time}가 높은 정확도를 보고한 것과 본 결과의 대비는 바로 이 해상도 차이로 설명된다.

\section{설계 결정에 대한 논의}
\label{sec:disc-design}

\paragraph{왜 킬만으로 교전을 정의하는가.} 근접성만으로 교전을 정의하면 서로 가까이 있으나 싸우지 않는 상황(파밍, 대치)이 대량으로 섞인다. 킬은 공개 텔레메트리에서 전투가 실제로 일어났음을 보장하는 유일한 사건이다. 대가는 킬 없는 교전의 누락이며, 이는 정의의 범위 제한으로 명시하고 검증 프로토콜에서 별도로 계량한다.

\paragraph{왜 위치를 선택적으로 0으로 두는가.} 지도 위치는 그래프 모델에 필수적인 구조 정보이지만, 순차 모델에서는 ``바론 둥지 근처의 교전은 특정 팀이 이긴다''와 같은 위치 편향을 암기하게 만든다. 따라서 $X^{\mathrm{node}}$에는 보존하고 거시 시퀀스에서는 0으로 둔다.

\paragraph{왜 스칼라 상태를 보간하지 않는가.} 아이템 구매, 버프 만료, 레벨 상승은 이산 전이이다. 프레임 사이를 선형 보간하면 실제로 존재하지 않는 매끄러운 중간 상태가 생기고, 특히 온셋 이후 프레임을 참조하는 보간은 교전 결과 자체를 입력에 새게 한다. 계단 유지는 이 두 문제를 동시에 피한다.

\paragraph{선택 편향의 회피.} 교전의 존재는 킬로 조건화하지만, 라벨 창의 \emph{결과}에는 조건화하지 않는다. 예측 지평에 특정 신호가 있는 창만 남기면 $P(Y \mid X, \text{signal}) \ne P(Y \mid X)$가 되어 Berkson 형 편향이 생기므로, 조용한 라벨 창도 그대로 둔다.

\section{코드--논문 감사}
\label{sec:disc-audit}

학술 논문의 방법 기술과 실제 코드 사이의 불일치는 재현성의 주요 위협이다. 본 연구는 최종본 제출 전에 코드와 논문을 항목별로 대조하는 감사를 수행하였고, 발견된 모든 불일치를 논문의 기술에 맞게 코드를 고치는 방향으로 해소하였다. 표 \ref{tab:audit}에 그 내역을 요약한다.

\begin{table}[htbp]
  \centering
  \caption{사전 감사에서 수정된 코드--논문 불일치}
  \label{tab:audit}
  \small
  \begin{tabularx}{\linewidth}{@{}lXX@{}}
    \toprule
    항목 & 발견 & 수정 \\
    \midrule
    A3 지름 분할 & 단일 연결(4{,}000 단위)로 분할하여 사슬형 군집의 지름이 상한을 초과할 수 있었음 & 완전 연결·지름 상한 분할로 교체 \\
    A1 온셋 검증 & ``생존 2 명(지도 전체)''과 ``반경 내 2 명(생존 무관)''을 별개 조건으로 검사 & 생존 $\wedge$ 반경 내를 단일 논리곱으로 검사 \\
    A2 인접 병합 & 15 초 / 2{,}000 단위 병합이 설정만 있고 실행되지 않음 & 병합 단계 구현 \\
    B1 특수 킬 & \texttt{CHAMPION\_SPECIAL\_KILL}을 두 번째 킬로 계수 & 보너스만 반영 \\
    MLP 경로 & 헤드라인 러너의 MLP가 95 차원 거시 시퀀스를 사용 & 2{,}980 차원 정합 입력으로 재경로 \\
    \bottomrule
  \end{tabularx}
\end{table}

수정의 순효과는 말뭉치가 1{,}115{,}123 건에서 994{,}365 건으로 11\% 줄어든 것(일방적 정리와 중복 재교전이 제거됨)과 LightGBM 시험 AUC가 $.675$에서 약 $.669$로 $0.006$ 낮아진 것이다(국소화 $-0.0017$, 라벨 $-0.0038$). 계열 간 순서(테이블 $\gg$ 신경망 표현)와 조건부 예측 가능성의 결론은 바뀌지 않았다. 본 논문은 심사를 거친 수치를 본문에 유지하고 이 절에서 수정 후 수치를 명시한다. \todo{학위논문 최종본에서는 수정된 코드로 전체 실험(3 시드 $\times$ 2 단계)을 재실행하여 표 \ref{tab:baseline}--\ref{tab:conditional}을 갱신하는 것을 권장한다.}

\section{타당성 위협}
\label{sec:disc-threats}

\paragraph{내적 타당성.} 무누출 계약은 단위 시험으로 강제되지만, Riot API 자체의 결함(예: 누락된 프레임, 지연된 사건 타임스탬프)은 통제할 수 없다. 교전의 3.67\%는 서로 다른 장소에서 동시에 벌어져 시간적으로 겹치는데, 독립 표본을 가정하는 모델에는 무해하지만 경기 전체를 순차 처리하는 모델에는 라벨 누출이 될 수 있어 클리핑·마스킹이 필요하다.

\paragraph{외적 타당성.} 데이터는 한국 서버 상위 티어 랭크 경기이다. 다른 지역, 낮은 티어, 프로 경기로의 일반화는 검증되지 않았다. 세 패치에 걸친 시간순 홀드아웃은 단기 공변량 이동만을 다루며, 시즌 단위의 큰 변화(지도 개편, 신규 오브젝트)는 별도 검증이 필요하다.

\paragraph{구성 타당성.} 교환가치 라벨의 가중치(표 \ref{tab:label-weights})는 도메인 지식에 근거한 설계 선택이지 데이터로부터 추정된 값이 아니다. 가중치를 바꾸면 라벨이 바뀌므로 결과는 이 정의에 상대적이다. 검출기가 찾은 교전이 사람의 판단과 일치하는지는 진행 중인 주석 연구의 결과에 달려 있다.

\section{한계}
\label{sec:disc-limitations}

(i) 리플레이 없이 공개 API만 사용하므로 초 단위 이하의 정보는 원천적으로 없다. (ii) 챔피언 스킬 세트와 조합 시너지는 임베딩을 통해서만 간접적으로 표현된다. (iii) 25 종 이상의 모델 $\times$ 3 시드 $\times$ 10 처치의 전수 실험은 계산 비용이 크다. (iv) 킬 없는 교전은 정의상 제외된다. (v) 라벨 가중치와 국소화 상수는 민감도 분석 없이 고정되었다.
